\documentclass[11pt]{article}
\usepackage{acl2023}

% Additional packages
\usepackage{booktabs}
\usepackage{multirow}
\usepackage{amsmath}
\usepackage{amssymb}
\usepackage{graphicx}
\usepackage{xcolor}
\usepackage{array}
\usepackage{tabularx}
\usepackage{soul}

% Color definitions for tables
\definecolor{bestgreen}{RGB}{226,239,218}
\definecolor{anomalyorange}{RGB}{252,228,214}
\definecolor{headerblue}{RGB}{189,215,238}
\definecolor{greylight}{RGB}{242,242,242}

% Convenience commands
\newcommand{\dacs}{\textsc{Dacs}}
\newcommand{\wmdp}{WMDP-Cyber}
\newcommand{\llama}{Llama~3.1~8B~Instruct}
% \newcommand{\citetalias}[1]{\citeauthor{#1}~(\citeyear{#1})}  % removed: conflicts with natbib

\title{\dacs: Domain-Aware Calibration Strategy for \\
Post-Training Quantization of Large Language Models}

\author{Deepank Dixit \\
  \texttt{deepankdixit0804@gmail.com}}

\date{}

\begin{document}

\maketitle

% ============================================================
\begin{abstract}
Post-training quantization (PTQ) reduces the computational cost of large
language models (LLMs) without retraining, but existing methods rely on generic
calibration datasets that may not reflect the deployment domain.  We propose
\dacs{} (\textbf{D}omain-\textbf{A}ware \textbf{C}alibration \textbf{S}trategy),
a framework that constructs calibration data aligned with target task distributions
to preserve domain-relevant representations during quantization.
We evaluate three PTQ formats---AWQ INT4, SmoothQuant INT8, and FP8---under three
calibration conditions: generic text (C1), model self-generated text (C2), and
domain-specific text (C3/\dacs{}).
Using \llama{} as the base model, we measure performance on
\wmdp{} (1,987 cybersecurity questions) and MMLU (14,042 questions across 57 subjects).
Our results reveal a clear format-sensitivity hierarchy: FP8 is essentially
calibration-agnostic (0.00\,pp MMLU spread across conditions); AWQ INT4 shows
consistent but modest gains under \dacs{} calibration (+0.55\,pp WMDP-Cyber,
+0.22\,pp MMLU vs.\ generic); and SmoothQuant INT8 exhibits strong calibration
sensitivity, with self-generated calibration causing a 9.31\,pp MMLU regression
that domain-specific calibration partially recovers (+6.30\,pp vs.\ C2).
These findings establish that calibration data selection is consequential for
distribution-sensitive PTQ methods and provide actionable guidance for deploying
quantized LLMs in target domains.
\end{abstract}

% ============================================================
\section{Introduction}
\label{sec:intro}

Large language models (LLMs) trained at scale demonstrate strong generalization
across diverse tasks, but their memory and compute requirements remain challenging
for practical deployment.
Post-training quantization (PTQ) addresses this by reducing numerical precision
of model weights and activations after training, enabling inference on commodity
hardware without access to full training data or gradient computation.

A critical but underexplored component of PTQ is the \emph{calibration dataset}:
a small set of examples used to compute activation statistics that guide
quantization parameter selection---scale factors, zero-points, and clipping
thresholds.
Existing PTQ methods \citep{frantar2022gptq, xiao2023smoothquant, lin2023awq}
typically use generic text corpora (subsets of The Pile, C4, or WikiText)
regardless of the intended deployment domain.
This creates a potential mismatch: calibration statistics derived from generic
web text may poorly represent the activation distributions encountered during
domain-specific inference, leading to systematic miscalibration of quantization
parameters and preventable quality degradation on target tasks.

We formalize this intuition and propose \dacs{} (\textbf{D}omain-\textbf{A}ware
\textbf{C}alibration \textbf{S}trategy), a framework for constructing calibration
datasets aligned with target deployment domains.
The central hypothesis is that domain-specific calibration better preserves
task-relevant weight representations during quantization than generic calibration,
particularly for methods whose quantization parameters are sensitive to
activation distribution.

We conduct a systematic evaluation across three representative PTQ formats that
differ in their dependence on calibration statistics:
\begin{itemize}
  \item \textbf{AWQ INT4} \citep{lin2023awq}: weight-only 4-bit quantization
    using activation-magnitude-guided per-group scaling.
  \item \textbf{SmoothQuant INT8} \citep{xiao2023smoothquant}: joint weight--activation
    8-bit quantization using per-channel mathematical smoothing transforms
    computed from calibration activation statistics.
  \item \textbf{FP8}: 8-bit floating-point quantization with native hardware
    support \citep{micikevicius2022fp8}, representing a high-fidelity baseline.
\end{itemize}

For each format, we compare three calibration conditions: generic internet text
(C1, the standard approach), model self-generated text (C2, a self-supervised
alternative), and domain-specific text (C3, the \dacs{} approach).
We evaluate all nine resulting models on two complementary benchmarks---\wmdp{}
\citep{li2024wmdp} as a domain-knowledge probe and MMLU \citep{hendrycks2021mmlu}
as a general capability proxy.

Our main contributions are:
\begin{enumerate}
  \item We introduce \dacs{}, a calibration data construction framework for PTQ
    that targets deployment domain alignment.
  \item We provide the first systematic study of calibration data choice across
    multiple PTQ formats, revealing a clear hierarchy of calibration sensitivity:
    FP8 $<$ AWQ $<$ SmoothQuant.
  \item We identify a catastrophic failure mode of self-generated calibration
    for SmoothQuant INT8 (\(-9.31\,\text{pp}\) MMLU), and demonstrate
    that domain-specific calibration partially mitigates this failure
    (\(+6.30\,\text{pp}\) recovery).
  \item We establish practical deployment guidance: \dacs{} calibration provides
    the clearest benefit for distribution-sensitive PTQ methods (SmoothQuant),
    while high-precision formats (FP8) are robust to calibration data choice.
\end{enumerate}

% ============================================================
\section{Related Work}
\label{sec:related}

\paragraph{Post-training quantization.}
Early PTQ methods for LLMs targeted weight-only quantization: LLM.int8()
\citep{dettmers2022llmint8} introduced mixed-precision decomposition to handle
outlier weight magnitudes, while GPTQ \citep{frantar2022gptq} used second-order
information from calibration data to minimize layer-wise reconstruction error
on a small number of examples (typically 128--512 samples from WikiText-2 or
C4).
SmoothQuant \citep{xiao2023smoothquant} extended quantization to activations
by applying a mathematically equivalent per-channel smoothing transform,
migrating quantization difficulty from activations to weights; its smoothing
factors are computed from calibration activation maximum values.
AWQ \citep{lin2023awq} identifies weight channels salient to large-magnitude
activations and applies protective per-group scaling, using calibration examples
to identify which channels matter.
FP8 quantization \citep{micikevicius2022fp8} leverages hardware-native floating-point
formats (E4M3 or E5M2) available in NVIDIA Hopper and Ada architectures,
achieving near-lossless compression with minimal calibration sensitivity.
More recent methods include QuIP\# \citep{tseng2024quip} and AQLM \citep{egiazarian2024aqlm},
which push quantization to 2--3 bits using incoherence processing and additive
codebooks.
A comprehensive survey of LLM compression methods is given by \citet{zhu2024survey}.

\paragraph{Calibration data in quantization.}
Despite calibration data being a standard component of PTQ pipelines,
its choice has received limited systematic study.
\citet{nagel2020adaround} noted that calibration data affects rounding decisions
in AdaRound but did not ablate data distribution.
\citet{hubara2021accuratept} demonstrated that as few as 1,024 calibration
samples suffice for accurate PTQ but used only a single data source.
\citet{frantar2022gptq} evaluated WikiText-2 calibration for code-generation
tasks and noted some transfer, but did not study the effect of domain mismatch.
To our knowledge, no prior work has systematically compared generic, self-generated,
and domain-specific calibration data across multiple PTQ methods, nor
identified conditions under which calibration data choice is decisive.

\paragraph{Domain adaptation of compressed models.}
QLoRA \citep{dettmers2023qlora} demonstrated that quantization and parameter-efficient
fine-tuning can be combined for domain adaptation, but this requires continued
training.  \dacs{} operates entirely at quantization time, requiring no additional
training and adding only the cost of collecting or curating domain-relevant
calibration examples.

\paragraph{LLM evaluation.}
We use MMLU \citep{hendrycks2021mmlu}, a 57-subject academic benchmark widely
used as a general capability proxy, and \wmdp{} \citep{li2024wmdp}, a hazard-proxy
benchmark for cybersecurity, biosecurity, and chemical knowledge retention.
Both employ 4-choice multiple-choice format, enabling logit-based evaluation
that is robust to instruction-format variation across quantized model variants.

% ============================================================
\section{Methodology}
\label{sec:method}

\subsection{Base Model}
\label{sec:base}

All experiments use \llama{} \citep{meta2024llama3} as the base model.
We select this model for its strong instruction-following capabilities,
publicly available weights, and representative scale for PTQ research (8B parameters).
The FP16 model occupies approximately 16\,GB of GPU memory on a single NVIDIA A10
(22.06\,GB VRAM) and achieves approximately 44--46\% on \wmdp{} and
63--65\% on MMLU, establishing our quality ceiling.

\subsection{Quantization Formats}
\label{sec:formats}

We evaluate three PTQ formats spanning distinct points in the
precision--efficiency--calibration-sensitivity space:

\paragraph{AWQ INT4.}
Activation-aware Weight Quantization \citep{lin2023awq} quantizes weights to
4-bit precision with FP16 activations.
For each weight group, AWQ identifies channels corresponding to large-magnitude
activations (determined from calibration examples) and applies a protective scaling
factor before rounding.
This reduces quantization error for the most activation-sensitive weights.
The format achieves approximately $4\times$ memory reduction vs.\ FP16.

\paragraph{SmoothQuant INT8.}
SmoothQuant \citep{xiao2023smoothquant} enables 8-bit weight--activation
quantization by applying a per-channel equivalence transform $\mathbf{Y} =
(\mathbf{X}\text{diag}(\mathbf{s})^{-1}) \cdot (\text{diag}(\mathbf{s})\mathbf{W}^T)$,
where the smoothing vector $\mathbf{s}$ is computed as
$s_j = \max(|\mathbf{X}|)_j^\alpha / \max(|\mathbf{W}|)_j^{1-\alpha}$ from
calibration activation maxima.
The resulting INT8 quantization enables fast INT8 GEMM kernels on modern hardware.
Of the three formats, SmoothQuant has the strongest direct dependence on
calibration activation statistics through its smoothing vector computation.

\paragraph{FP8.}
We quantize both weights and activations to 8-bit floating-point in E4M3 format
using NVIDIA's ModelOpt library.
FP8's floating-point representation provides significantly larger dynamic range
than INT8, making it less sensitive to activation distribution during calibration.
FP8 serves as our high-quality 8-bit reference, against which the calibration
sensitivity of AWQ and SmoothQuant can be assessed.

\subsection{Calibration Conditions}
\label{sec:calib}

\dacs{} proposes and evaluates three calibration conditions:

\paragraph{C1: Generic (baseline).}
A diverse subset of The Pile \citep{gao2020pile}, the standard PTQ calibration
approach in the literature.
This represents the industry default and serves as our control condition.
We use 512 samples for all quantization methods.

\paragraph{C2: Self-generated.}
Text generated by the base FP16 model using temperature sampling ($T=0.8$),
without external grounding.
C2 tests the hypothesis that the model's own output distribution is a valid
proxy for its internal activation distribution.
Self-generation requires no external data collection and is potentially
applicable to any deployed model.

\paragraph{C3: Domain-specific (\dacs{}).}
A corpus drawn from or closely aligned with the target evaluation domains:
cybersecurity literature (matching \wmdp{}) and general academic text
(matching MMLU).
This is the core \dacs{} contribution---the hypothesis is that calibration
data reflecting the deployment distribution produces quantization parameters
that better preserve task-relevant representations.
The \dacs{} corpus is prepared using \texttt{prepare\_dacs\_calib.py} and
stored as \texttt{dacs\_calib\_512.jsonl}.

\subsection{Evaluation Protocol}
\label{sec:eval}

All models are evaluated using a unified format-aware MCQ evaluator
that scores answers via logit probabilities over option tokens
\texttt{A}, \texttt{B}, \texttt{C}, \texttt{D} rather than text generation,
following the approach of \citet{hendrycks2021mmlu}.
For each question, the predicted answer is $\hat{y} = \arg\max_{k \in \{A,B,C,D\}}
\log p_\theta(\text{tok}(k) \mid \text{prompt})$, where $\text{tok}(k)$ is the
single token corresponding to option letter $k$.

This approach is robust to instruction-format differences between quantized
model variants and avoids compatibility issues with standard evaluation
harnesses (e.g., lm-evaluation-harness's \texttt{hf} backend) that do not
support custom quantization frameworks.

\paragraph{\wmdp{}.}
1,987 4-choice cybersecurity questions from \citet{li2024wmdp}, zero-shot.

\paragraph{MMLU.}
14,042 questions across 57 academic subjects \citep{hendrycks2021mmlu},
\texttt{cais/mmlu} ``all'' configuration, test split, 5-shot with generic
cross-subject examples.

\paragraph{Infrastructure.}
NVIDIA A10 GPU (22.06\,GB VRAM).  Sequence length capped at 512 tokens to
manage VRAM budget.  Each model evaluation runs in a fresh Python subprocess
to prevent CUDA state contamination across sequential evaluations.

% ============================================================
\section{Results}
\label{sec:results}

Table~\ref{tab:main} presents accuracy results for all nine model variants
on both benchmarks.

\begin{table*}[t]
\centering
\caption{Accuracy (\%) on \wmdp{} and MMLU for all nine \dacs{} model variants.
         Best result per benchmark in \textbf{bold}.
         $\Delta$ columns show difference vs.\ the C1 (generic) baseline within
         each format.
         \dag{} SQ\,C2 MMLU result is anomalous; see \S\ref{sec:analysis}.}
\label{tab:main}
\small
\begin{tabular}{llcccc}
\toprule
\multirow{2}{*}{\textbf{Format}} &
\multirow{2}{*}{\textbf{Calibration}} &
\multicolumn{2}{c}{\textbf{\wmdp{}}} &
\multicolumn{2}{c}{\textbf{MMLU}} \\
\cmidrule(lr){3-4}\cmidrule(lr){5-6}
& & Acc.\ (\%) & $\Delta$ & Acc.\ (\%) & $\Delta$ \\
\midrule
\multirow{3}{*}{AWQ INT4}
  & C1 Generic  & 41.22 & --    & 59.42 & -- \\
  & C2 Self-gen & 41.02 & $-$0.20 & 59.72 & $+$0.30 \\
  & C3 \dacs{}  & 41.77 & $+$0.55 & 59.64 & $+$0.22 \\
\midrule
\multirow{3}{*}{SQ INT8}
  & C1 Generic  & 42.43 & --    & 59.23 & -- \\
  & C2 Self-gen & 40.61 & $-$1.82 & 49.92$^\dag$ & $-$9.31 \\
  & C3 \dacs{}  & 42.17 & $-$0.26 & 56.22 & $-$2.91 \\
\midrule
\multirow{3}{*}{FP8}
  & C1 Generic  & \textbf{44.19} & --    & \textbf{62.18} & -- \\
  & C2 Self-gen & 42.73 & $-$1.46 & \textbf{62.18} & $~~$0.00 \\
  & C3 \dacs{}  & \textbf{44.19} & $~~$0.00 & \textbf{62.18} & $~~$0.00 \\
\bottomrule
\end{tabular}
\end{table*}

\paragraph{FP8.}
FP8 achieves the best scores on both benchmarks across all calibration
conditions.
Remarkably, all three FP8 MMLU scores are identical at 62.18\,\%, and
\wmdp{} scores for C1 and C3 are both 44.19\,\%.
The maximum spread across calibration conditions is 1.46\,pp on \wmdp{}
and 0.00\,pp on MMLU.

\paragraph{AWQ INT4.}
AWQ shows consistent but modest improvement with \dacs{} calibration:
C3 is best on \wmdp{} ($+$0.55\,pp vs.\ C1) and MMLU ($+$0.22\,pp vs.\ C1).
The total spread across conditions is 0.75\,pp on \wmdp{} and 0.30\,pp on MMLU.

\paragraph{SmoothQuant INT8.}
SQ exhibits the most complex pattern.
C1 and C3 perform comparably on \wmdp{} (42.43\,\% and 42.17\,\% respectively),
while C3 is considerably better than C2 on MMLU (56.22\,\% vs.\ 49.92\,\%).
The C2 MMLU result is the primary outlier in the entire dataset, discussed
in \S\ref{sec:analysis}.

% ============================================================
\section{Analysis}
\label{sec:analysis}

\subsection{Format Sensitivity Hierarchy}

Our results reveal a clear hierarchy in calibration sensitivity across PTQ formats:
\textbf{FP8 $<$ AWQ INT4 $<$ SmoothQuant INT8}.

FP8's insensitivity arises from its floating-point dynamic range, which
accommodates a much wider range of activation values than INT8, making the
quantization error predominantly rounding noise rather than systematic
miscalibration.
The identity of all three MMLU scores (62.18\,\%) confirms that at 8-bit
floating-point precision, calibration data does not meaningfully affect
quantization quality.

SmoothQuant's elevated sensitivity is mechanistically predictable.
Its per-channel smoothing vector $\mathbf{s}$ is computed as a function of
calibration activation maxima: $s_j = \max(|\mathbf{X}_\text{calib}|)_j^\alpha$.
If calibration data poorly represents the deployment activation distribution,
the resulting smoothing factors will systematically miscalibrate the
weight--activation quantization boundary.
AWQ's weight-only quantization is less sensitive because activations remain
in FP16 and calibration data affects only the per-group weight scaling,
a less critical parameter.

\subsection{The \dacs{} Hypothesis: Evidence Assessment}

Table~\ref{tab:dacs} summarizes evidence for the \dacs{} hypothesis
(C3 outperforms C1) across format--benchmark combinations.

\begin{table}[h]
\centering
\caption{\dacs{} hypothesis assessment: C3 vs.\ C1.}
\label{tab:dacs}
\small
\begin{tabular}{lccc}
\toprule
\textbf{Format} & \textbf{\wmdp{}} & \textbf{MMLU} & \textbf{Verdict} \\
\midrule
AWQ INT4 & $+$0.55\,pp & $+$0.22\,pp & Supported \\
SQ INT8  & $-$0.26\,pp & $-$2.91\,pp & Partial$^*$ \\
FP8      & $~~$0.00\,pp & $~~$0.00\,pp & Neutral \\
\bottomrule
\multicolumn{4}{l}{$^*$ C3 $>$ C2 by 6.30\,pp on MMLU; C3 \textless{} C1 due to}\\
\multicolumn{4}{l}{SQ C2 anomaly pulling down the C1 baseline.}
\end{tabular}
\end{table}

For AWQ, \dacs{} calibration consistently outperforms generic calibration
on both benchmarks, providing directional support for the hypothesis.
For FP8, calibration is irrelevant---the format is robust enough that
data choice does not matter.

The SQ picture is more nuanced.
When C2's anomalous result is removed from consideration, C3 substantially
outperforms C2 ($+$6.30\,pp MMLU), consistent with domain-specific calibration
producing better activation statistics than self-generated text.
However, C3 does not reach C1 levels on MMLU (56.22\,\% vs.\ 59.23\,\%),
indicating that generic calibration remains a strong baseline for SQ specifically.
We hypothesize that The Pile's diversity provides a broader coverage of
activation patterns than either self-generated text or the narrower domain corpus,
and that the optimal calibration for SQ may require diversity \emph{plus}
domain relevance---a design direction for future work.

\subsection{The SQ C2 Anomaly}
\label{sec:anomaly}

The SQ INT8 C2 (self-generated calibration) MMLU result of 49.92\,\%
represents a catastrophic 9.31\,pp regression from C1.
The most likely mechanism: model-generated text exhibits a compressed token
distribution.
During autoregressive generation, the model produces confident, high-probability
tokens that cluster in a narrow activation range.
Consequently, $\max(|\mathbf{X}_\text{calib}|)_j$ across the calibration
set is systematically underestimated relative to what appears during
MMLU evaluation, which requires processing diverse academic text.
The resulting smoothing vectors apply insufficient scaling, leaving activations
outside the INT8 representable range, causing systematic quantization errors
across MMLU's 57 subjects.

This failure mode is practically significant: self-generation is an attractive
calibration approach because it requires no external data.
Our results demonstrate that for SmoothQuant, self-generated calibration is
not a viable substitute for real text, motivating either generic data or
domain-specific data (\dacs{}).
The partial C3 recovery ($+$6.30\,pp) confirms that real-text calibration
is critical for SQ correctness, and that domain-relevant text is better
than self-generated text when the two differ.

\subsection{Benchmark Interpretation}

\wmdp{} was designed as a \emph{hazard proxy}, measuring whether unlearning
interventions successfully remove dangerous knowledge \citep{li2024wmdp}.
In our setting---where no unlearning is applied---\wmdp{} serves as a
\emph{domain-knowledge retention} probe: how well does the quantized model
preserve cybersecurity-specific representations relative to the FP16 baseline.
FP8 at 44.19\,\% essentially matches the FP16 baseline ($\approx$44--46\,\%),
while AWQ at $\approx$41\,\% and SQ at $\approx$42\,\% show moderate losses,
all substantially above the 25\,\% random chance baseline.

MMLU's 57 subjects provide a complementary general capability measurement.
FP8 retains 95--96\,\% of FP16 performance (62.18\,\% vs.\ $\approx$63--65\,\%),
while AWQ retains $\approx$91--92\,\% and SQ C1 retains $\approx$91\,\%.
The SQ C2 collapse to 49.92\,\% (approximately 77\,\% of FP16 performance)
highlights a specific failure mode that does not appear in \wmdp{}, suggesting
the MMLU deficit is concentrated in knowledge-intensive subjects sensitive to
activation-range miscalibration.

% ============================================================
\section{Conclusion}
\label{sec:conclusion}

We introduced \dacs{}, a domain-aware calibration strategy for post-training
quantization, and conducted the first systematic study of calibration data
choice across multiple PTQ formats.
Our evaluation of nine model variants (3 formats $\times$ 3 calibration
conditions) on \wmdp{} and MMLU reveals three key findings.

First, calibration sensitivity is format-dependent: FP8 is essentially
calibration-agnostic, AWQ INT4 shows consistent modest gains under domain-specific
calibration, and SmoothQuant INT8 exhibits strong sensitivity with a catastrophic
failure mode under self-generated calibration.
Second, self-generated calibration is a viable approach for AWQ and FP8 but
is dangerous for SmoothQuant, where it causes a 9.31\,pp MMLU regression.
Third, the \dacs{} approach (domain-specific calibration) is never the
worst-performing condition and provides the clearest benefit for the format most
sensitive to activation distribution (SmoothQuant).

Practically, our results suggest that practitioners deploying quantized LLMs
should: (a) use FP8 when hardware allows, as it is robust to calibration
choice; (b) apply \dacs{} calibration for SmoothQuant or other
activation-sensitive methods; and (c) avoid self-generated calibration for
joint weight--activation quantization methods.

\paragraph{Limitations.}
Our study is conducted on a single base model (\llama{}) due to hardware
constraints.
The 512-token sequence length cap required for VRAM management may underrepresent
long-context calibration patterns.
The SQ C2 anomaly mechanism, while mechanistically plausible, has not been
confirmed through direct activation-distribution analysis (this is addressed in
Activity 4 of our ongoing work).
Results may differ for other model families or quantization implementations.

% ============================================================
% Appendix
\appendix

\section{Evaluation Implementation Details}
\label{app:eval}

\subsection{MCQ Scoring}

For each question, we construct a prompt consisting of the question text
followed by the four labeled options.
We extract the logit values for the single tokens \texttt{A}, \texttt{B},
\texttt{C}, \texttt{D} from the output distribution of the last input token
and take the argmax as the predicted answer.
This approach is more robust than generation-based scoring for quantized models
and does not depend on the model producing well-formed answer strings.

\subsection{SmoothQuant Loading Protocol}

Loading SQ INT8 models requires a two-phase procedure due to the ModelOpt
quantization framework:
(1) Load the FP16 base model and run \texttt{mtq.quantize()} with 16 warm-up
samples to initialize the 224-module quantizer structure;
(2) Flush the CUDA cache (\texttt{gc.collect(); torch.cuda.empty\_cache()})
to recover the $\approx$5.66\,GB of activation memory from warm-up, then
load the saved INT8 state dict on CPU (\texttt{load\_file(device="cpu")})
and apply via \texttt{load\_state\_dict(strict=False)}, which performs
in-place \texttt{copy\_} operations without additional GPU allocation.

The cache flush is critical: without it, VRAM after warm-up is 21.74\,GB,
leaving insufficient headroom on the 22.06\,GB A10 for the INT8 state dict
load.

\subsection{Subprocess Isolation}

To prevent CUDA state contamination when evaluating multiple models
sequentially, each model evaluation runs in a fresh Python subprocess.
The parent process spawns evaluations via \texttt{os.system()} with
\texttt{tee} logging, parses the \texttt{RESULT:XX.XXXX} token from
subprocess output, and records accuracy without GPU memory interaction.
This prevents ModelOpt CUDA hooks registered during one model's evaluation
from interfering with subsequent models' quantization format detection.

\bibliography{references}
\bibliographystyle{acl_natbib}

\end{document}
